% Timing diagram of a job: the first jobs of a picture, following job module
% 6 from the scheduler to the start of its iterator. Rendered to
% timing_job.png by "make diagrams" (see the DIAGRAMS variable in the
% top-level Makefile). It is shown in ALGORITHM.md, in the section "Timing of
% a job and a result" of the Dispatcher, next to the table of the clock cycles
% of a job, and the clock cycles t, t+1, ... are the same as in the table.
%
% All values are read off a simulation of the first instance of the
% dispatcher in sim/dispatcher_tb.vhd (i_dispatcher: 16 job modules in groups
% of 5, so job module 6 is in group 1), in the first picture: t is the clock
% cycle from 325 ns to 335 ns (the rising edge of clk at 325 ns). Anything that
% changes the processes or the registers of the dispatcher, or that
% testbench, has to change this file too.
%
% Drawing convention, the same as in timing_result.tex: single bits are always
% drawn with their true value. A bus is white in the clock cycles where it
% carries the job of job module 6, light grey with its value where it carries
% another job, and grey without a value where it carries nothing (its valid
% flag is low). The values of cx are in hex, as the 18 bits of the
% register.
%
% The scheduler is started by start_i in t-1 (sched_active_r is high from t,
% and its delayed copy active_d from t+1, not drawn), so job_idx_valid_o is
% high from t+2. All the job modules are idle, so after that the scheduler
% selects one job module in each clock cycle, and job_cx_r moves on by stepx
% (0x333 here) for each job. The busy flags in grp_busy_r(1) are those of the
% job module 5+(cnt_r mod 5) of group 1 in the clock cycle before: job
% module 6 in t+1 and in t+6, and job module 7 in t+7.

\documentclass[border=8pt]{standalone}

\usepackage[T1]{fontenc}
\usepackage{helvet}
\renewcommand{\familydefault}{\sfdefault}

\usepackage{tikz-timing}

% White: the job of job module 6. Light grey: another job. Grey (U): nothing.
\tikzset{
   timing/d/background/.style={fill=white},
   timing/u/background/.style={fill=black!25},
   wh/.style={timing/d/background/.style={fill=white}},
   gr/.style={timing/d/background/.style={fill=black!8}},
   timing/name/.style={font=\ttfamily\small},
   timing/d/text/.style={font=\scriptsize},
}

\begin{document}
\begin{tikztimingtable}[timing/slope=0.08, timing/coldist=6pt, xscale=2.2, yscale=1.4]
clk                                  & 10{HL} \\
start\_i                             & 2L 2H 16L \\
job\_scheduler/cnt\_r                & [gr] 2D{4} 2D{5} [wh] 2D{6} [gr] 2D{7} 2D{8} 2D{9} 2D{10} 2D{11} 2D{12} 2D{13} \\
job\_busy\_i(6)                      & 14L 6H \\
job\_scheduler/grp\_busy\_r(1)       & 16L 4H \\
job\_idx\_valid\_o                   & 8L 12H \\
job\_idx\_start\_o                   & 8U [wh] 2D{6} [gr] 2D{7} 2D{8} 2D{9} 2D{10} 2D{11} \\
job\_start\_r(6)                     & 10L 2H 8L \\
job\_started\_r, tbl\_valid\_o       & 10L 10H \\
tbl\_idx\_o                          & 10U [wh] 2D{6} [gr] 2D{7} 2D{8} 2D{9} 2D{10} \\
job\_cx\_r                           & 4U [wh] 8D{30000} [gr] 2D{30333} 2D{30666} 2D{30999} 2D{30CCC} \\
job\_start\_o(6)                     & 12L 2H 6L \\
job\_cx\_o(1)                        & 6U [wh] 8D{30000} [gr] 2D{30333} 2D{30666} 2D{30999} \\
job\_started\_d                      & 12L 8H \\
job(6)/res\_start\_r                 & 14L 2H 4L \\
job(6)/res\_cx\_r                    & 14U [wh] 6D{30000} \\
job\_started\_dd                     & 14L 6H \\
grp\_job\_busy\_r(1)                 & 16L 4H \\
\extracode
\begin{pgfonlayer}{background}
   \vertlines[densely dotted, black!50]{0,2,...,20}
\end{pgfonlayer}
\foreach \x/\l in {0/t-2, 2/t-1, 4/t, 6/t+1, 8/t+2, 10/t+3, 12/t+4, 14/t+5, 16/t+6, 18/t+7}
   \node[above, font=\small] at (\x+1, 0.9) {\l};
\end{tikztimingtable}
\end{document}
